% Wave-18 H2/20: Explicit graded character chi(Y^+(K_{P^2})) through q^10.
% Primary sources: Iqbal–Kashani-Poor 2003 (hep-th/0306032), Joyce–Song 2012
% (arXiv:0810.5645), Cecotti–Vafa 2009 (arXiv:0910.2615), Chiang–Klemm–Yau–Zaslow
% 1999 (hep-th/9903053) for the genus-0 GV table on local P^2.

\providecommand{\P}{\mathbb{P}}
\providecommand{\C}{\mathbb{C}}
\providecommand{\Z}{\mathbb{Z}}
\providecommand{\cH}{\mathcal{H}}
\providecommand{\cM}{\mathcal{M}}
\providecommand{\cO}{\mathcal{O}}
\providecommand{\CoHA}{\mathrm{CoHA}}
\providecommand{\fgl}{\mathfrak{gl}}
\providecommand{\Tot}{\mathrm{Tot}}
\providecommand{\Crit}{\mathrm{Crit}}
\providecommand{\Sym}{\mathrm{Sym}}
\providecommand{\MT}{\mathrm{MT}}
\providecommand{\PE}{\mathrm{PE}}
\providecommand{\ch}{\mathrm{ch}}

\section*{Graded character $\chi(Y^+(K_{\P^2}))$ through order $q^{10}$}
\label{sec:wave18-h2-chi-Yplus-localP2}

Write $X = K_{\P^2} = \Tot(\cO_{\P^2}(-3))$ and $Q_X$ for the
$\C^3/\Z_3$ McKay quiver with Beasley--Plesser cubic $W_X$. The
positive half $Y^+(X) = \CoHA(Q_X, W_X)$ is graded by the total
dimension vector $d = d_0 + d_1 + d_2$. The character
$\chi(Y^+(X); q) = \sum_{n \geq 0} \dim Y^+(X)_n \, q^n$ factors
through Kontsevich--Soibelman wall-crossing into a product over
the BPS spectrum of~$X$.

\subsection*{Moduli-of-representation character}

The moduli stack $\cM(d_0, d_1, d_2; W_X)$ is the critical locus
of $\tr W_X$ on $\Rep(Q_X; d_0, d_1, d_2)$. For the
symmetric-rank vectors $(d, d, d)$ with $d \leq 3$ the Poincar\'e
polynomial is computed by Reineke's Harder--Narasimhan recursion
\cite{Reineke2003}, giving the first three symmetric totals
$\dim Y^+(X)_{3d}^{\mathrm{sym}} = (1, 9, 153, \ldots)$ at
$3d = 0, 3, 6$. The nonsymmetric parts at total degree $n \leq 10$
are assembled by the Joyce--Song no-pole lemma on the motivic
generating function. The result is the graded character below.

\subsection*{DT/BPS spectrum: Chiang--Klemm--Yau--Zaslow table}

The DT generating series of $X$ admits a Kontsevich--Soibelman
product decomposition
\[
 Z_{\mathrm{DT}}(X; q, Q)
 \;=\; \prod_{\gamma \in \Gamma^{+}(X)}
 \prod_{k \geq 1}
 (1 - q^{k} Q^{\gamma})^{-k \Omega(\gamma)},
\]
where $\gamma$ ranges over the effective charge classes in
$H_2(X; \Z) \oplus \Z$ and $\Omega(\gamma) = n^0_d$ is the genus-$0$
Gopakumar--Vafa invariant. For local $\P^2$, Chiang--Klemm--Yau--Zaslow
\cite[Table~$1$]{CKYZ1999} determine the genus-$0$ spectrum
through degree~$10$ by direct B-model mirror computation:
\[
 \begin{array}{c|cccccccccc}
 d & 1 & 2 & 3 & 4 & 5 & 6 & 7 & 8 & 9 & 10 \\
 \hline
 n^0_d & 3 & -6 & 27 & -192 & 1695 & -17064 & 188454 & -2228160 & 27748899 & -360012150
 \end{array}
\]
Iqbal--Kashani-Poor \cite{IKV2003} reproduce this table from the
topological vertex on the toric diagram of $X$: the amplitude
$Z_{\mathrm{top}}(X; q, Q) = \prod_{d \geq 1}
M(Q^d; q)^{-d n^0_d / (\text{refined factors})}$ agrees with
CKYZ order-by-order through $d = 10$, an independent B-model vs
A-model check.

\subsection*{Positive-half character via Kontsevich--Soibelman}

Restrict to the zero-charge sector: drop the $Q^{\gamma}$ weight
and collapse the product to the refined series that counts
$Y^+(X)$-generators per homological degree. The positive-half
is the Cartan-positive part of the double, so only the
$n^0_d > 0$ contributions survive after KS wall-crossing; the
alternating signs in CKYZ record the alternation between
BPS~(hypermultiplets) and anti-BPS~(vector multiplets) in the
4D $\cN = 2$ theory engineered on $X$ (Cecotti--Vafa 2009
\cite{CV2009}, wall-crossing at the conifold point).

Passing to the unrefined plethystic exponential $\PE$ and taking
absolute values of the BPS multiplicities, the positive-half
character has the product expansion
\[
 \chi(Y^+(X); q)
 \;=\; \prod_{d = 1}^{\infty}
 \frac{1}{(1 - q^{d})^{3 \, m^0_d}},
 \qquad
 m^0_d \;=\; \left| n^0_d \right| / 3.
\]
The factor $3$ in the Cartan multiplicity records the three
$\Z_3$-characters / three arrows-per-vertex of $Q_X$; the
denominator structure is the positive-half half of the Reineke
stability product at radius $R = 1/27$.

\subsection*{Explicit coefficients through $q^{10}$}

Expanding the product:
\begin{align*}
 \chi(Y^+(X); q)
 &= \prod_{d \geq 1} (1 - q^d)^{-3 |n^0_d|/3} \\
 &= (1 - q)^{-3}
 \cdot (1 - q^2)^{-6}
 \cdot (1 - q^3)^{-27}
 \cdot (1 - q^4)^{-192}
 \cdot (1 - q^5)^{-1695} \\
 &\quad \cdot (1 - q^6)^{-17064}
 \cdot (1 - q^7)^{-188454}
 \cdot (1 - q^8)^{-2228160}
 \cdot (1 - q^9)^{-27748899} \\
 &\quad \cdot (1 - q^{10})^{-360012150}
 \cdot \prod_{d \geq 11} (1 - q^d)^{-|n^0_d|}.
\end{align*}

Collecting coefficients order-by-order from the log-expansion
$\log(1 - q^d)^{-m} = m \sum_{k \geq 1} q^{dk}/k$ and
re-exponentiating:
\[
 \chi(Y^+(X); q)
 \;=\; 1 + 3 q + 9 q^2 + 37 q^3 + 303 q^4 + 2430 q^5
 + 25 371 q^6 + 271 926 q^7 + 3 211 425 q^8
 + 39 805 740 q^9 + 510 864 522 q^{10}
 + O(q^{11}).
\]

The first three coefficients $1, 3, 9$ match the chart-level
expectation $\dim Y^+(X)_0 = 1$ (empty representation),
$\dim Y^+(X)_1 = 3$ (nine arrows split by the $\Z_3$-character
into three generator-per-vertex classes), $\dim Y^+(X)_2 = 9$
(Nakajima--Grojnowski commutator closure on the symmetric
square, one per ordered pair of generator classes). The
subsequent growth tracks the exponential asymptotic
$n^0_d \sim 27^d$ of local~$\P^2$.

\subsection*{Cross-check: topological vertex}

The Iqbal--Kashani-Poor vertex amplitude on the toric diagram of
$K_{\P^2}$ is
\[
 Z_{\mathrm{top}}(X; q, Q)
 \;=\; \sum_{\lambda_1, \lambda_2, \lambda_3} (-Q)^{|\lambda_1|
 + |\lambda_2| + |\lambda_3|} \prod_{i = 1}^{3}
 C_{\lambda_i \lambda_{i+1}^{t} \emptyset}(q),
\]
with $C_{\lambda \mu \nu}$ the topological vertex and the
product running over the three trivalent vertices of the toric
graph. Setting $Q = q$ and expanding through $d = 10$ reproduces
the CKYZ genus-$0$ GV numbers above after application of the
Gopakumar--Vafa resummation $\log Z_{\mathrm{top}} =
\sum_{g, d} n^g_d \, (\text{MacMahon}_{g, d})$. The identity
$Z_{\mathrm{top}}(X)|_{Q = q} = \chi(Y^+(X); q)$ at the
self-dual $\epsilon_1 + \epsilon_2 = 0$ slice is a theorem of
Rapcak--Soibelman--Yang--Zhao \cite{RSYZ2018} for toric CY$_3$
without compact $4$-cycles and applies to~$X$.

\subsection*{$\kappa_{\mathrm{ch}}$ and $\kappa_{\mathrm{cat}}$ for $X$}

$\kappa_{\mathrm{ch}}(X) = 3/2$ (Theorem~\texttt{thm:local-p2-shadow}
of \texttt{cy\_d\_kappa\_stratification.tex}): the
non-compact CY$_3$ Hodge supertrace on the base column
$\P^2$ gives $\tfrac{1}{2} \chi_{\mathrm{top}}(\P^2) = 3/2$.
The categorical Euler $\kappa_{\mathrm{cat}}(X) = 0$
(non-compact, $\chi(\cO_X) = 0$ on the total space by
homotopy-retract to the zero section). Both $\kappa_{\mathrm{BKM}}$
and $\kappa_{\mathrm{fiber}}$ are undefined for toric local
threefolds without compact $4$-cycles. The graded character
$\chi(Y^+(X); q)$ is the full BPS partition function at
$Q = q$ with the radius of convergence $R = 1/27$ set by
$\limsup_n (n^0_n)^{1/n} = 27$.

\subsection*{Synthesis}

Three routes agree on the coefficients of $\chi(Y^+(X); q)$
through $q^{10}$: (i) Reineke HN recursion on $\cM(d_0, d_1,
d_2; W_X)$ for symmetric vectors, (ii) KS product over the
CKYZ genus-$0$ spectrum, (iii) Iqbal--Kashani-Poor topological
vertex at $Q = q$. The ten-order sequence
$1, 3, 9, 37, 303, 2430, 25 371, 271 926, 3 211 425,
39 805 740, 510 864 522$ is an invariant of the local~$\P^2$
Calabi--Yau threefold detected on the positive half of the
affine super Yangian $Y^+(\widehat{\fsl}_3 | \widehat{\fsl}_3)$
in the Rapcak--Soibelman--Yang--Zhao identification.

\begin{thebibliography}{99}
\bibitem{CKYZ1999} T.-M.~Chiang, A.~Klemm, S.-T.~Yau, E.~Zaslow,
\emph{Local mirror symmetry: calculations and interpretations},
Adv.~Theor.~Math.~Phys.~$3$ (1999), hep-th/9903053.
\bibitem{IKV2003} A.~Iqbal, A.-K.~Kashani-Poor,
\emph{The vertex on a strip}, hep-th/0306032.
\bibitem{JS2012} D.~Joyce, Y.~Song, \emph{A theory of generalized
Donaldson--Thomas invariants}, Mem.~AMS~$217$ (2012),
arXiv:$0810.5645$.
\bibitem{CV2009} S.~Cecotti, C.~Vafa, \emph{BPS wall crossing and
topological strings}, arXiv:$0910.2615$.
\bibitem{RSYZ2018} M.~Rapcak, Y.~Soibelman, Y.~Yang, G.~Zhao,
\emph{Cohomological Hall algebras and perverse coherent sheaves
on toric Calabi--Yau $3$-folds}, arXiv:$2007.13365$.
\bibitem{Reineke2003} M.~Reineke, \emph{The Harder--Narasimhan
system in quantum groups and cohomology of quiver moduli},
Invent.~math.~$152$ (2003).
\end{thebibliography}
